\قسمت{تغییر مسیر}

شما قصد دارید یک کوتاه‌کننده لینک را پیاده‌سازی کنید. کوتاه‌کننده لینک برای لینک‌های طولانی لینک‌های کوتاه و قابل به خاطر سپردن را در اختیار کاربران قرار می‌دهد.
برای پیاده‌سازی نیاز دارید تا کاربر با فراخوانی یکی از \متن‌لاتین{endpoint}های شما و دادن لینک کوتاه به لینک اصلی \متن‌لاتین{redirect} شود.
در نظر داشته باشید که این \متن‌لاتین{redirect} باید به گونه‌ای باشد که کاربر آدرس جدید را \متن‌لاتین{cache} نکند.
برای این امر از کدام \متن‌لاتین{HTTP Status Code} می‌بایست استفاده کرد؟

\نمره{۲}

% پاسخ پیشنهادی است و پیش از استفاده در تصحیح نیاز به بازبینی دارد.
\begin{پاسخ}

کد \متن‌لاتین{302 Found} (یا \متن‌لاتین{307 Temporary Redirect} که روش تقاضا را هم حفظ می‌کند).

این‌ها تغییر مسیر موقت هستند و مرورگر نشانی تازه را به جای نشانی کوتاه نمی‌نشاند؛
در حالی که \متن‌لاتین{301} و \متن‌لاتین{308} دائمی‌اند و مرورگر آن‌ها را نگه می‌دارد،
که برای کوتاه‌کننده‌ی لینک یعنی از دست دادن امکان تغییر مقصد و از دست دادن آمار.

\شروع{فقرات}
\فقره نام بردن \متن‌لاتین{302} یا \متن‌لاتین{307}: ۷۵ درصد نمره
\فقره بیان اینکه \متن‌لاتین{301} به دلیل \متن‌لاتین{cache} شدن مناسب نیست: ۲۵ درصد نمره
\پایان{فقرات}

\end{پاسخ}
